\section{Conclusion}

This project aimed to develop an operating system inspired by the Apollo Guidance
Computer (AGC), focusing on its scheduling and error-handling algorithms. The result,
Apollo-RTOS, reimagines the core ideas of the AGC with modern concepts, demonstrating
that the AGC's algorithms were ahead of their time and remain relevant for efficient
operating system design.

The strength of these algorithms lies in their simplicity, a theme central to
Apollo-RTOS's implementation. This simplicity makes the system easy to understand,
explain, and verify compared to more complex designs. Apollo-RTOS shows that even with
the \texttt{micro:bit}'s limited resources—a 16MHz CPU and 16KB of RAM—careful design
can achieve significant functionality. The system can concurrently run 16 processes,
serialize their execution deterministically, and handle errors gracefully, making it a
strong candidate for real-world applications.

However, there is still room for improvement. Testing in this project was limited, and
it was impossible to explore every possible combination of events within the given
timeframe. Many hardware features were not tested under adverse conditions, which is
crucial before real-world deployment. Additionally, the implementation is not fully
optimized. Memory and CPU usage were managed leniently, without strict restrictions on
execution time. For an RTOS, strict guarantees on upper bounds for each operation are
essential, and Apollo-RTOS currently lacks these guarantees or has not yet been tested
for them.

In terms of features, Apollo-RTOS does not yet support all hardware components
available on the \texttt{micro:bit}. Notably, it lacks drivers for the SPI bus,
watchdog timer, sensors, radio device, and others. Some limitations also stem from the
hardware constraints of the \texttt{micro:bit} v1. For example, the Cortex M0 processor
lacks a floating-point unit (FPU), requiring software libraries to handle
floating-point arithmetic. Integrating such a library remains a task for future work.

Despite these limitations, Apollo-RTOS highlights how historical concepts can still be
applied in modern contexts. It also demonstrates the feasibility of creating a general
purpose, lightweight, efficient, and reliable operating system for resource-constrained
microcontrollers, opening new possibilities for real-world applications.

